%-----------------------------------------------------------------
% Rodrigo Chang, 2019
% Banco de Guatemala
% Programa de Estudios Superiores 2019-2020
% Diapositivas del curso "Programación I"
%
% Descargo de responsabilidad.
%   El contenido de este documento es responsabilidad exclusiva de
% su(s) autores. Cualquier opinión expresada no necesariamente
% representa la postura institucional del Banco de Guatemala. 
% 
% Contenido.
%   El contenido de estas diapositivas está construido a partir de 
% modificaciones al material provisto por Ana Bell, Eric Grimson y 
% John Guttag para el curso 6.0001 Introduction to Computer Science 
% and Programming in Python, otoño 2016, Massachusetts Institute of 
% Technology: MIT OpenCourseWare, https://ocw.mit.edu
%
% Licencia: 
% Creative Commons BY-NC-SA
% https://creativecommons.org/licenses/by-nc-sa/4.0/.
%
% Formato. 
%   El formato utilizado para la elaboración de este documento 
% ha sido modificado a partir del repositorio: 
%       https://github.com/rinnert/statlearning-python-pwf
%
% Kurt Rinnert, 2019
%   Copying and distribution of this file, with or without 
% modification, are permitted in any medium without royalty provided 
% the copyright notice and this notice are preserved. This file is 
% offered as-is, without any warranty.
%
%-----------------------------------------------------------------
\documentclass[mathserif]{beamer}
% Para configurar las diapositivas en relación 16:9, 
% agregar la opción aspectratio=169

% Título de la presentación y tópicos a tratar
\usepackage{arev, arevmath}
\usepackage[scaled]{cabin}
%\usepackage[T1]{fontenc}
\usepackage[super]{nth}
\usepackage{pifont}
\usepackage{wasysym}
\usepackage{tabularx}
\usepackage{array}
\usepackage{booktabs}
\usepackage{boldline}
\usepackage{colortbl}
\usepackage{amsmath}
\usepackage{bm}
\usepackage{tcolorbox}
\usepackage{adjustbox}
\usepackage{minibox}
\usepackage{makecell}
\usepackage{adjustbox}
\usepackage{textcomp}
\usepackage[absolute,overlay]{textpos}
\setlength{\TPHorizModule}{1mm}%
\setlength{\TPVertModule}{1mm}%
\tcbuselibrary{skins}
\usepackage{graphicx}
\usepackage[spanish, es-tabla]{babel}
\usepackage[utf8]{inputenc}
\usepackage{csquotes}

\usepackage{comment}
\usepackage{svg}

\usepackage{verbatim}
\usepackage{pdfpages}

% Columnas múltiples
\usepackage{multicol}

% -------------------------------------------
%	LISTINGS
% -------------------------------------------
% Paquetes necesarios para los códigos
\usepackage{listings}
\usepackage{color}
% Configuración de listings (pendiente)
%\input{conf/listingStyles}
% Redefinición de titulos de caption
\renewcommand\lstlistingname{\textbf{Código}}
\renewcommand\lstlistlistingname{Códigos}

% Definición de colores para Python (Jupyter notebooks)
\definecolor{python_keyword_color}{RGB}{0, 128, 0}
\definecolor{python_string_color}{RGB}{186, 34, 34}
\definecolor{python_comment_color}{RGB}{1, 121, 121}
\definecolor{frame_numbers_color}{rgb}{0.5,0.5,0.5}
\definecolor{python_identifier_color}{RGB}{0,0,0}
\definecolor{jupyter_background}{RGB}{247, 247, 247}
\definecolor{jupyter_frame}{RGB}{207, 207, 207}

% Definicion de estilo para Python
\lstdefinestyle{custom_python}{
	language=python,			% the language of the code
	backgroundcolor=\color{jupyter_background}, 
	% the size of the fonts that are used for the code
	basicstyle=\linespread{1}\small\ttfamily,	
	% Configuración de colores
	commentstyle=\color{python_comment_color},	% comment style
	identifierstyle=\color{python_identifier_color},
	keywordstyle=\mdseries\color{python_keyword_color},	% keyword style
	numberstyle=\tiny\color{frame_numbers_color},	
	stringstyle=\color{python_string_color}, % string literal style 
	captionpos=t,				% sets the caption-position to bottom
	breakatwhitespace=false,	% sets if automatic breaks should only happen at whitespace
	breaklines=true,			% sets automatic line breaking
	morekeywords={},			% if you want to add more keywords to the set
	deletekeywords={eps},		% delete keywords from the given language
	escapeinside={\%*}{*)},		% if you want to add LaTeX within your code
	framesep=3pt,
	keepspaces=true,			% keeps spaces in text,       
	numbers=none, 	% (none, left, right)
	numbersep=5pt,	% how far the line-numbers are from the code
	% if not set, the frame-color may be changed on line-breaks within not-black text (e.g. comments (green here))
	rulecolor=\color{jupyter_frame}, 
	% show spaces everywhere adding particular underscores; it overrides 'showstringspaces'
	showspaces=false,                
	% underline spaces within strings only
	showstringspaces=false,          
	% show tabs within strings adding particular underscores
	showtabs=false,                  
	stepnumber=0,	% the step between two line-numbers      
	% sets default tabsize to 2 spaces
	tabsize=2,                       
	% lets you use non-ASCII characters; for 8-bits encodings only, does not work with UTF-8
	extendedchars=true,      
	% Support additional characters            
	literate=
  {á}{{\'a}}1 {é}{{\'e}}1 {í}{{\'i}}1 {ó}{{\'o}}1 {ú}{{\'u}}1
  {Á}{{\'A}}1 {É}{{\'E}}1 {Í}{{\'I}}1 {Ó}{{\'O}}1 {Ú}{{\'U}}1
  {à}{{\`a}}1 {è}{{\`e}}1 {ì}{{\`i}}1 {ò}{{\`o}}1 {ù}{{\`u}}1
  {À}{{\`A}}1 {È}{{\'E}}1 {Ì}{{\`I}}1 {Ò}{{\`O}}1 {Ù}{{\`U}}1
  {ä}{{\"a}}1 {ë}{{\"e}}1 {ï}{{\"i}}1 {ö}{{\"o}}1 {ü}{{\"u}}1
  {Ä}{{\"A}}1 {Ë}{{\"E}}1 {Ï}{{\"I}}1 {Ö}{{\"O}}1 {Ü}{{\"U}}1
  {â}{{\^a}}1 {ê}{{\^e}}1 {î}{{\^i}}1 {ô}{{\^o}}1 {û}{{\^u}}1
  {Â}{{\^A}}1 {Ê}{{\^E}}1 {Î}{{\^I}}1 {Ô}{{\^O}}1 {Û}{{\^U}}1
  {œ}{{\oe}}1 {Œ}{{\OE}}1 {æ}{{\ae}}1 {Æ}{{\AE}}1 {ß}{{\ss}}1
  {ű}{{\H{u}}}1 {Ű}{{\H{U}}}1 {ő}{{\H{o}}}1 {Ő}{{\H{O}}}1
  {ç}{{\c c}}1 {Ç}{{\c C}}1 {ø}{{\o}}1 {å}{{\r a}}1 {Å}{{\r A}}1
  {€}{{\EUR}}1 {£}{{\pounds}}1,   
}

% Configurar el listing para el documento actual
\lstset{
	style = custom_python, 
	frame = single
}

% -------------------------------------------
%	OTRAS CONFIGURACIONES
% -------------------------------------------
\makeatletter
\newcommand{\antsize}{\@setfontsize{\antsize}{4pt}{4pt}}
\makeatother
\newcommand{\at}{\makeatletter @\makeatother}

\newcommand{\cmark}{\ding{51}}%
\newcommand{\bottomline}[1]{\vskip0pt plus 1fill{\alert{#1}\phantom{g}\vskip 1.0mm}}

% Rutas para las imágenes
%\graphicspath{{figs/}{figs/common/}}

% Separación de filas en las tablas más amplia
\renewcommand{\arraystretch}{1.2}


% -------------------------------------------
%	ESTILO DE LA PRESENTACIÓN
% -------------------------------------------

\mode<presentation>%
{
	\usetheme{default}
	%\usetheme[width=2.5cm]{PaloAlto}
	\usecolortheme{dove}
	\useoutertheme{infolines}
	% oder auch nicht

	% BANGUAT Colors
	\definecolor{BANGUATblue}{RGB}{37,86,162} % 0x255682
	\definecolor{BANGUATorange}{RGB}{255,130,0} % 0xff8200
	\definecolor{BANGUATgreen}{RGB}{0,100,0}
	\definecolor{BANGUATdark}{RGB}{80,80,80} % 0x505050
	\definecolor{BANGUATlight}{RGB}{120,120,120}
	\definecolor{BANGUATbrown}{RGB}{178,91,0}

	\definecolor{codebg}{rgb}{0.95,0.95,0.95}

	% Color theme
	\setbeamercolor{alerted text}{fg=BANGUATorange}
	\setbeamercolor{frametitle}{fg=BANGUATblue}
	\setbeamercolor{title}{fg=BANGUATblue}
	\setbeamercolor{subtitle}{fg=BANGUATorange}
	\setbeamercolor{normal text}{fg=BANGUATdark}
	\setbeamercolor{author in foot}{fg=BANGUATblue}
	\setbeamercolor{item}{fg=BANGUATblue}
	\setbeamercolor{footline}{fg=BANGUATblue}
	%\setbeamercolor{item projected}{bg=BANGUATorange}
	%\setbeamercolor{item projected}{fg=white}

	\setbeamertemplate{headline}
	{}
	\setbeamertemplate{frametitle}
	{
		%\textbf{{\insertframetitle\phantom{g}}}\\
		%\textbf{\insertframetitle\phantom{g}}\\
		\textbf{\underline{\insertframetitle\phantom{g}}}\\
		%\textbf{\underline{\insertframetitle}}\\
		\vskip 1.0mm
		%{\color{UOLgold}\hrule height 2pt}
		%\par
	}
	\addtobeamertemplate{frametitle}{}{\vspace{-1em}}
	\setbeamertemplate{footline}{
		{%
			\textbf{ \hskip 3.0mm\insertshorttitle\phantom{.}---\phantom{.}\insertshortinstitute\hfill\insertframenumber\,/\,\inserttotalframenumber\hskip 3.0mm} 
		}
	}

	\setbeamertemplate{navigation symbols}{}%remove navigation symbols
	\setbeamertemplate{itemize items}[circle]
	\setbeamertemplate{enumerate items}[fg=BANGUATblue]
	\setbeamercolor{itemize items}{fg=BANGUATblue}
	\setbeamercolor{sidebar}{bg=BANGUATblue}
	\setbeamercolor{title in sidebar}{fg=BANGUATorange}
	\setbeamercolor{author in sidebar}{fg=BANGUATorange}
	\setbeamercolor{section in sidebar}{fg=BANGUATorange}
}



% -------------------------------------------
%	CONFIGURACIÓN DE DIAPOSITIVA DE TÍTULO
% -------------------------------------------

% Título del curso
\title[Programación I]{\textbf{%
Introducción a las ciencias de la computación\\%
{\it y programación en Python}}\\%
%\href{https://ocw.mit.edu/courses/electrical-engineering-and-computer-science/6-0001-introduction-to-computer-science-and-programming-in-python-fall-2016/}{\tiny\it Basado en el curso ``Introduction to Computer Science and Programming in Python'' por Ana Bell, Eric Grimson, y John Guttag.}%
\vspace{3em}%
}
%\vspace{-2.5cm}{}

% Autor y correo electrónico
\author{\href{mailto:?to=Rodrigo Chang <rrcp@banguat.gob.gt>&subject=PES Programación I}{Rodrigo Chang}}

% Información del instituto
\institute[\href{https://www.banguat.gob.gt/}{Banco de Guatemala}]%
{\color{BANGUATblue} \bfseries \href{https://www.banguat.gob.gt/}{Banco de Guatemala}\\%
\vspace{3mm}%
\includegraphics[width=0.13\textwidth]{figs/common/logo-banguat.png}}

% Sin fecha en particular
\date{}

% Información de licencias y copyright
\titlegraphic{
	 %\texorpdfstring{\vspace{-3cm}}{}
	 \vspace{-1.5cm}
	 \begin{minipage}[b][1cm][b]{0.25\textwidth}%
	    \color{BANGUATlight}\antsize
		%Copyright \textcopyright~2019\\
		\href{mailto:?to=Rodrigo Chang <rrcp@banguat.gob.gt>&subject=PES Programación I}{Rodrigo Chang <rrcp@banguat.gob.gt>}\\
		Este material está construido a partir de modificaciones al material provisto por Ana Bell, Eric Grimson y John Guttag para el curso \textit{6.0001 Introduction to Computer Science and Programming in Python}, otoño 2016, Massachusetts Institute of Technology: MIT OpenCourseWare, \href{https://ocw.mit.edu}{https://ocw.mit.edu}. Licencia: \href{https://creativecommons.org/licenses/by-nc-sa/4.0/}{Creative Commons BY-NC-SA}.
		
		%Copying and distribution of this file, with or without modification, are permitted in any medium without royalty provided the copyright notice and this notice are preserved.  This file is offered as-is, without any warranty.


		%Some of the figures in this presentation are taken from ``An Introduction to Statistical Learning, with applications in R''  (Springer, 2013) with permission from the authors: G. James, D. Witten,  T. Hastie and R. Tibshirani 
	 \end{minipage}%
	 \hspace{9cm}
}

% Restar uno al contador
\addtocounter{framenumber}{-1}

% Cajas de colores
\newcommand{\tabboxset}{\tcbset{enhanced, nobeforeafter, boxrule=0pt, boxsep=0pt, colback=codebg, colframe=codebg, coltext=BANGUATdark, rounded corners, arc=4pt, fonttitle={\bfseries\tiny}}}
\newcommand{\codeboxset}{\tcbset{enhanced, nobeforeafter, boxrule=0pt, boxsep=0pt, colback=codebg, colframe=codebg, coltext=BANGUATdark, rounded corners, arc=4pt, fonttitle={\bfseries\tiny}}}

% Comandos útiles
\newcommand{\orange}{\color{BANGUATorange}}
\newcommand{\blue}{\color{BANGUATblue}}
\newcommand{\dark}{\color{BANGUATdark}}
\newcommand{\R}{\mathbb{R}}
\newcommand{\dat}[1]{{\footnotesize\tt\orange #1}}
\newcommand{\mdat}[1]{\text{\dat{#1}}}
\newcommand{\val}[1]{{\footnotesize\tt\blue #1}}
\newcommand{\mval}[1]{\text{\val{#1}}}
\newcommand{\e}[1]{\emph{#1}}
\newcommand{\bh}{\hat{\beta}}
\newcommand{\h}{\hat}
\newcommand{\bl}{\bottomline}


\makeatletter
\newcommand{\includegraphicsdpi}[3]{%
	\pdfimageresolution=#1%
	\includegraphics[#2]{#3}%
	\pdfimageresolution=72%
}

\newcommand{\Quote}[2]{%
	\begin{center} 
		\begin{minipage}{0.7\textwidth} 
			\hrule
			\vskip 3mm
			\emph{{\color{BANGUATblue} ``#1''}\\
			{\color{BANGUATorange} --- #2}}
			\vskip 3mm
			\hrule
			\vskip 2mm
		\end{minipage}
	\end{center}
}

% Texto centrado en la diapositiva, quizás para el abstract o para indicar abajo de un quote
\newenvironment{blurb}%
	{\begin{center}\begin{minipage}{0.6\textwidth}\footnotesize}
	{\end{minipage}\end{center}}

% Environment para códigos o algoritmos. Parecido a un bottomline, se puede utilizar para dar definiciones o enfatizar en algún concepto
\newenvironment{cpage}%
	{\begin{center}\begin{minipage}{0.75\textwidth}}
	{\end{minipage}\end{center}}

% Bloque (similar a un pop-up) de alerta para indicar resultados en tablas, definiciones, fórmulas, etc.
% Primer argumento es el tamaño 
% Segundo argumento es el título del bloque
\newenvironment{popblock}[2]%
	{\begin{center}\begin{minipage}{#1}\footnotesize
		\begin{tcolorbox}[colframe=codebg, colback=white, colupper=BANGUATdark, title={\normalsize\bfseries\blue #2}]}
	{\end{tcolorbox}\end{minipage}\end{center}}
\makeatother

% Incluir diapositivas en PDF
% 1 - Título de la diapositiva
% 2 - páginas del PDF a incluir
% 3 - ruta hacia el archivo PDF
\newcommand{\includeslides}[3]{%
{%
\setbeamercolor{background canvas}{bg=}%
\addtocounter{framenumber}{1}%
\includepdf[pages={#1}, link=true, pagecommand={%
\Large\blue\bfseries\underline{%
#3\phantom{g}}%
\vspace{-1em}%
\vskip 1.0mm}]{#2}%
}}
\subtitle{%
  \bfseries{Programación orientada a objetos}\\%
  %{\tiny\it Manejo de strings, \textit{guess and check}, métodos de aproximación}\\%
}

% Inicio del documento
\begin{document}
\frame[plain]{
	%\vskip 1.0mm
	\titlepage
	%\vskip 1.0mm
}
%-----------------------------------------------------------------
% Empieza el contenido de la presentación
%-----------------------------------------------------------------

\begin{frame}{Abstract}
	\Quote{Commenting your code is like cleaning your bathroom — you never want to do it, but it really does create a more pleasant experience for you and your guests.}{Ryan Campbell}
    
	\begin{blurb}
		Veremos conceptos básicos y ejemplos del paradigma de programación orientada a objetos.
	\end{blurb}
\end{frame}


\section{Programación orientada a objetos}

% Objetos
\begin{frame}[fragile]{Objetos}
    \begin{itemize}
        \item Como vimos antes, Python representa información de diferentes formas:
\begin{lstlisting}
1234    3.14159     "Hello"
[2, 4, 6, 8, 10]
{"GT": "Guatemala", "HN":"Honduras"}
\end{lstlisting}
        
        \item Cada uno es un {\blue\bfseries objeto}, el cual tiene:
        \begin{itemize}
            \item Un {\orange tipo}.
            \item Una {\orange representación interna} (primitiva o compuesta).
            \item Un conjunto de procedimientos de {\orange interacción} con el objeto.
        \end{itemize}
        
        \item Un objeto es una {\blue\bfseries instancia} de un tipo:
        \begin{itemize}
            \item \lstinline|1234| es una instancia de un entero :  \lstinline|int|.
            \item \lstinline|"hello"| es una instancia de un \textit{string}: \lstinline|str|.
        \end{itemize}
    \end{itemize}
\end{frame}

% Programación orientada a objetos
\begin{frame}[fragile]{Programación orientada a objetos}
    \begin{itemize}
        \item En Python, {\blue\bfseries ¡TODO ES UN OBJETO!}
        
        \item Podemos {\orange\bfseries crear nuevos objetos} de algún tipo.
        
        \item Podemos {\orange\bfseries manipular objetos}.
        \item Podemos {\orange\bfseries destruir objetos}.
        \begin{itemize}
            \item Utilizando \lstinline|del| o ``dejándolos'' a un lado\footnote{Python utiliza un sistema de recolección de memoria para objetos destruidos o inaccesibles, llamado ``recolección de basura''.}.
        \end{itemize}
    \end{itemize}
\end{frame}

% Objetos nuevamente - representación interna
\begin{frame}[fragile]{Ya en serio: ¿qué son los objetos?}
    \begin{itemize}
        \item Representan una {\orange\bfseries abstracción de datos} que captura:
        
        \item Una {\blue\bfseries representación interna}:
        \begin{itemize}
            \item A través de \emph{atributos}.
        \end{itemize}
        
        \item Una {\blue\bfseries interfaz} para interactuar con el objeto:
        \begin{itemize}
            \item A través de \emph{métodos} ({\footnotesize alias procedimientos/funciones})
            \item Que definen el comportamiento, pero {\blue ocultan detalles de implementación}.
        \end{itemize}
    \end{itemize}
\end{frame}

% Ejemplo con listas
\includeslides{5}{slides-60001/MIT6_0001F16_Lec8_mod.pdf}{Ejemplo: \lstinline|[1,2,3,4]| es de tipo lista}

% Ventajas de la OOP
\begin{frame}[fragile]{Ventajas de la OOP}
    \begin{itemize}
        \item Permiten {\blue\bfseries empaquetar} datos y procedimientos para trabajar sobre estos con interfaces bien definidas.
        
        \item Desarrollo al estilo ``{\blue\bfseries divide y conquistarás}''
        \begin{itemize}
            \item Implementación y pruebas sobre cada clase separada.
            \item Incrementan la modularidad y reducen la complejidad.
        \end{itemize}
        
        \item Permiten {\blue\bfseries reutilizar} el código fácilmente.
        \begin{itemize}
            \item Muchos módulos definen nuevas clases.
            \item Cada clase tiene un ambiente separado (no hay problemas de nombres de funciones).
            \item La {\orange herencia} permite redefinir o extender el comportamiento de una clase padre. 
        \end{itemize}
    \end{itemize}
\end{frame}

% Crear y usar un tipo propio
\begin{frame}[fragile]{Las clases implementan tipos propios}
    \begin{itemize}
        \item Existe una diferencia entre {\blue\bfseries crear una clase} e {\blue\bfseries instanciar una clase}.
        
        \item {\blue\bfseries Crear} una clase involucra
        \begin{itemize}
            \item Definir el nombre de la clase.
            \item Definir los atributos.
            \item Por ejemplo: \textit{alguien definió la clase \lstinline|list|, sus atributos y métodos}.
        \end{itemize}
        
        \item {\blue\bfseries Utilizar} la clase involucra:
        \begin{itemize}
            \item Crear una nueva {\orange instancia} del objeto.
            \item Realizar operaciones con la instancia.
            \item Por ejemplo: \lstinline|L=[1,2]| y \lstinline|len(L)|.
        \end{itemize}
    \end{itemize}
\end{frame}


% Definición de una clase
\begin{frame}[fragile]{Definición de una clase}
    \begin{itemize}
        \item Para definir un nuevo tipo, utilizamos la palabra \lstinline|class|:
\begin{lstlisting}
class Coordinate(object):
    #define attributes here
\end{lstlisting}

         \item Similar a \lstinline|def|, {\blue indentamos} para indicar qué elementos pertenecen a la clase.
         
         \item \lstinline|object| se refiere a la {\blue clase padre} de \lstinline|Coordinate|.
         \begin{itemize}
             \item \lstinline|Coordinate| es una subclase de \lstinline|object|.
             \item \lstinline|object| es una superclase de \lstinline|Coordinate|.
         \end{itemize}
    \end{itemize}
\end{frame}

% ¿Qué son atributos?
\begin{frame}[fragile]{¿Qué son atributos?}
    \begin{itemize}
        \item Datos y funciones que {\orange\bfseries pertenecen} a la clase.
        \item {\orange\bfseries Atributos de datos}
        \begin{itemize}
            \item Objetos de datos que componen la clase.
            \item Por ejemplo: en \lstinline|Coordinate| un atributo es la coordenada $x$.
        \end{itemize}
        
        \item {\orange\bfseries Métodos}
        \begin{itemize}
            \item Funciones que solamente funcionan con esta clase.
            \item Permiten {\blue interactuar} con el objeto.
            \item En \lstinline|Coordinate| podríamos definir un método \lstinline|distance| que nos devuelva la distancia hacia otro objeto \lstinline|Coordinate|. 
            \begin{itemize}
                \item Notar que no necesariamente habría significado de \lstinline|distance| entre dos objetos \lstinline|list|.
            \end{itemize}
        \end{itemize}
    \end{itemize}
\end{frame}

% El método constructor
\begin{frame}[fragile]{El método constructor}
    \begin{itemize}
        \item Primero debemos definir {\orange\bfseries cómo crear} una instancia de la clase.
        
        \item A este método le llamamos {\orange\bfseries constructor}. En Python, definimos la función \lstinline|__init__(self, ...)|.
    \end{itemize}
    
\begin{lstlisting}
class Coordinate(object):
    def__init__(self, x, y):
        self.x = x
        self.y = y
\end{lstlisting}
\end{frame}

% Instanciando una clase
\begin{frame}[fragile]{Instanciando una clase}
\begin{lstlisting}
c = Coordinate(3,4)
origin = Coordinate(0,0)
print(c.x, origin.x)
\end{lstlisting}

    \begin{itemize}
        \item Los atributos de datos de una instancia se convierten en {\orange\bfseries variables de la instancia}.
        
        \item Notar que:
        \begin{itemize}
            \item No proveemos argumento \lstinline|self| $\Rightarrow$ automático.
            \item Utilización del \lstinline|.| para acceder a los atributos/métodos.
            \item Al llamar a \lstinline|Coordinate|, invocamos la función \lstinline|__init__|.
        \end{itemize}
    \end{itemize}
\end{frame}

% ¿Qué es un método?
\begin{frame}[fragile]{¿Qué es un método?}
    \begin{itemize}
        \item Es también un atributo, pero {\orange\bfseries solo funciona con esta clase}.
        
        \item Python siempre pasa el objeto como primer argumento:
        \begin{itemize}
            \item Utilizamos \lstinline|self| como el primer argumento en cualquier método.
        \end{itemize}
        
        \item {\blue\bfseries El operador \lstinline|.|} permite acceder a cualquier atributo. 
        \begin{itemize}
            \item Llamar a un método del objeto.
            \item Acceder a los datos almacenados en el objeto.
        \end{itemize}
    \end{itemize}
\end{frame}

% Nuestro primer método
\begin{frame}[fragile]{Nuestro primer método}
\begin{lstlisting}
class Coordinate(object):
    """ A coordinate made up of an x and y value """
    def __init__(self, x, y):
        """ Sets the x and y values """
        self.x = x
        self.y = y
    def distance(self, other):
        """ Returns the euclidean distance between two points """
        x_diff_sq = (self.x - other.x)**2
        y_diff_sq = (self.y - other.y)**2
        return (x_diff_sq + y_diff_sq)**0.5
\end{lstlisting}
    \begin{itemize}
        \item Aparte de \lstinline|self| y la notación \lstinline|.|, los métodos {\orange\bfseries se comportan igual que las funciones}.
        \begin{itemize}
            \item Toman parámetros.
            \item Realizan cómputos, procedimientos. 
            \item Devuelven valores.
        \end{itemize}
    \end{itemize}
    
    \bottomline{¿Qué asume el método \lstinline|distance|?}
\end{frame}

% Cómo utilizar un método
\begin{frame}[fragile]{Cómo utilizar un método}
    \begin{columns}
        \begin{column}{0.5\textwidth}
        \begin{itemize}
            \item Forma convencional: lo invocamos sobre un objeto.
\begin{lstlisting}
c = Coordinate(3,4)
zero = Coordinate(0,0)
print(c.distance(zero))
\end{lstlisting}
            \item Notar que se omite \lstinline|self|, pues está implicado por el objeto \lstinline|c|.
        \end{itemize}
        \end{column}
        \begin{column}{0.5\textwidth}
        \begin{itemize}
            \item Equivalente a:
\begin{lstlisting}
c = Coordinate(3,4)
zero = Coordinate(0,0)
print(Coordinate.distance(c, zero))
\end{lstlisting}
            \item Se operan los dos objetos entre sí invocando el método {\blue desde la clase} y no desde una instancia.
        \end{itemize}
        \end{column}
    \end{columns}
\end{frame}

% Imprimir la representación de un objeto
\begin{frame}[fragile]{}
\begin{lstlisting}
c = Coordinate(3,4)
print(c)
# Imprime: <__main__.Coordinate object at 0x7fa918510488>
\end{lstlisting}
    \begin{itemize}
        \item Representación {\blue\bfseries poco informativa} por defecto.
        
        \item Podemos definir un método \lstinline|__str__| para una clase.
        
        \item Cuando utilizamos \lstinline|print| sobre un objeto, Python llama a su método \lstinline|__str__|.
        
        \item Podemos escoger qué mostrar, supongamos que queremos:
    \end{itemize}
\begin{lstlisting}
print(c)
# Queremos: <3, 4>
\end{lstlisting}
\end{frame}

% Definición propia para print
\begin{frame}[fragile]{Definición propia para \lstinline|print|}
\begin{lstlisting}
class Coordinate(object):
    """ A coordinate made up of an x and y value """
    def __init__(self, x, y):
        """ Sets the x and y values """
        self.x = x
        self.y = y
    def distance(self, other):
        """ Returns the euclidean distance  """
        x_diff_sq = (self.x-other.x)**2
        y_diff_sq = (self.y-other.y)**2
        return (x_diff_sq + y_diff_sq)**0.5
    def __str__(self):
        """ Returns a string representation of self """
        return "<"+str(self.x)+","+str(self.y)+">"
\end{lstlisting}
    \bottomline{¿Qué tipo es devuelto por \lstinline|__str__|?}
\end{frame}

% Acerca de los tipos y las clases
\includeslides{17}{slides-60001/MIT6_0001F16_Lec8_mod.pdf}{Acerca de los tipos y las clases}

% Operadores especiales
\includeslides{18}{slides-60001/MIT6_0001F16_Lec8_mod.pdf}{Operadores especiales}

% Ejemplo de fracciones
\begin{frame}[fragile]{Ejemplo: fracciones}
    \begin{itemize}
        \item Crearemos un {\blue\bfseries nuevo tipo} para representar fracciones.
        
        \item Su {\blue\bfseries representación interna} serán dos enteros:
        \begin{itemize}
            \item Numerador.
            \item Denominador.
        \end{itemize}
        
         \item Definiremos una interfaz (métodos) que permitirán:
         \begin{itemize}
             \item Sumar, restar.
             \item Imprimir la fracción, convertir a flotante.
             \item Obtener el recíproco.
         \end{itemize}
    \end{itemize}
    \bottomline{Veamos el código}.
\end{frame}

% El poder de OOP
\begin{frame}[fragile]{El poder de OOP}
    \begin{itemize}
        \item Podemos {\blue\bfseries empaquetar objetos} que comparten:
        \begin{itemize}
            \item Atributos comunes.
            \item Procedimientos que operan sobre esos atributos.
        \end{itemize}
        
        \item Utilizamos {\blue\bfseries abstracción} para diferenciar cómo implementar un objeto y cómo utilizarlo.
        
        \item Podemos construir {\blue\bfseries capas} de abstracción de objetos que {\orange hereden} comportamiento de otras clases de objetos.
        
         \item Podemos crear {\blue\bfseries nuestras propias clases de objetos} utilizando las clases base de Python.
    \end{itemize}
\end{frame}

%-----------------------------------------------------------------
% Termina el contenido de la presentación
%-----------------------------------------------------------------

\begin{comment}

% Slide with two columns, half and half
\begin{frame}[fragile]{}
    \begin{columns}
        \begin{column}{0.5\textwidth}
        \end{column}
        \begin{column}{0.5\textwidth}
        \end{column}
    \end{columns}
\end{frame}

% Slide with bullets
\begin{frame}[fragile]{}
    \begin{itemize}
        \item 
    \end{itemize}
\end{frame}

% Slide with bullets and subbullets
\begin{frame}[fragile]{}
    \begin{itemize}
        \item {\blue\bfseries A}
        \begin{itemize}
            \item a
        \end{itemize}
    \end{itemize}
\end{frame}

% Slide with bullets and code
\begin{frame}[fragile]{}
    \begin{itemize}
        \item 
    \end{itemize}
\begin{lstlisting}
Code goes here
\end{lstlisting}
\end{frame}

% Slides PDF
\includeslides{54-57}{slides-60001/MIT6_0001F16_Lec6_mod.pdf}{Eficiencia: Fibonacci con diccionarios}

% Pytutor
\href{URL}{Ver en Python Tutor}


\end{comment}

\end{document}